\documentclass[exercice,a4]{thedocument}%
\usepackage[french]{babel}%
\usepackage{tabularx}%
\usepackage{array}%
\usepackage{enumitem}%
\startdatakeys
    \source{22GENMATAN1}
    \BO{2026}
    \cycle{4}
    \competencesDuSocle{CA,RA}
    \objectifsApprentissage{C4C2aa,C4C2ab,C4C2ac,C4D2ae}
    \motsCles{pourcentage,étendue,médiane,moyenne}
\enddatakeys
\begin{document}%

\medskip

\begin{tabularx}{\linewidth}{|X|}\hline
Pour être en bonne santé, il est recommandé d'avoir régulièrement une pratique physique. Une recommandation serait de faire au moins une heure de pratique physique par jour en moyenne. Sur 1,6 million d'adolescents de 11 à 17 ans interrogés, 81\,\% d'entre eux ne respectent pas cette recommandation.\\ \hline
\end{tabularx}
\begin{flushright}\small {\emph{D'après un communiqué de presse sur la santé}}\end{flushright}

\medskip

\begin{enumerate}
\item Sur les $1,6$ million d'adolescents de 11 à 17 ans interrogés, combien ne respectent pas cette recommandation ?
\end{enumerate}

\medskip

Après la lecture de ce communiqué, un adolescent se donne un objectif.

\begin{center} \textbf{Objectif: \og  \emph{Faire au moins une heure de pratique physique par jour en moyenne.} \fg}\end{center}

Pendant 14 jours consécutifs, il note dans le calendrier suivant, la durée quotidienne qu'il consacre à sa pratique physique:

\begin{center}
\begin{tabularx}{\linewidth}{|*{7}{>{\centering \arraybackslash}X|}}\hline
\textbf{Jour 1} 	&\textbf{Jour 2}&\textbf{Jour 3}		& \textbf{Jour 4}	&\textbf{Jour 5}	&\textbf{Jour} 6		&\textbf{Jour} 7\\ \hline
50 min	&15 min&1 h			&1 h 40 min	&30 min	&1 h 30 min	&40 min\\ \hline
\textbf{Jour 8}	&\textbf{Jour 9}&\textbf{Jour 10}		&\textbf{Jour 11}	&\textbf{Jour 12}&\textbf{Jour 13}	&\textbf{Jour 14}\\ \hline
15 min	&1 h	&1 h 30 min	&30 min		&1 h 	&1 h 		&0 min\\ \hline
\end{tabularx}
\end{center}

\begin{enumerate}[resume]
\item
	\begin{enumerate}
		\item Quelle est l'étendue des 14 durées quotidiennes notées dans le calendrier ?
		\item Donner une médiane de ces 14 durées quotidiennes.
	\end{enumerate}
\item
	\begin{enumerate}
		\item Montrer que, sur les 14 premiers jours, cet adolescent n'a pas atteint son objectif.
		\item Pendant les 7 jours suivants, cet adolescent décide alors de consacrer plus de temps au sport pour atteindre son objectif sur l'ensemble des $21$ jours.

Sur ces 7 derniers jours, quelle est la durée totale de pratique physique qu'il doit au minimum prévoir pour atteindre son objectif ?
	\end{enumerate}
\end{enumerate}

\bigskip

\showthesource
\end{document}
